% =============================================================================
% CAPÍTULO 4 — GUÍA OPERATIVA: PERFIL RECLUTADOR  (generado por tools/gen_ch4.py)
% =============================================================================
\chapter{Guía Operativa: Perfil Reclutador}
\label{ch:reclutador}
\resumenCapitulo{Este capítulo guía a la persona usuaria con perfil Reclutador en el
uso del panel de reclutamiento: la búsqueda inteligente y filtrada de talento, la
revisión y evaluación de portafolios, la gestión del pipeline de selección, la
comunicación directa con los candidatos y la configuración de su cuenta.}

\section{Visión General del Panel de Reclutamiento}
\label{sec:rec-dashboard}
Al iniciar sesión como Reclutador accederá al \campo{Dashboard}, que reúne los
indicadores clave de su actividad (perfiles vistos, guardados, chats y contrataciones),
un resumen automático de desempeño y gráficos de mercado. Desde la barra superior se
navega a las secciones \campo{Buscar Talento}, \campo{Rankings}, \campo{Pipelines} y
\campo{Chat}.

\captura{assets/img/ethoshub-reclutador/01_dashboard/DashboardReclutador1.png}{Para ir al apartado de «Dashboard» debe hacer clic al botón «Dashboard».}
\captura{assets/img/ethoshub-reclutador/01_dashboard/DashboardReclutador2.png}{Al observar el botón «Vistos» podemos apreciar el conteo de candidatos vistos en los últimos 7 días.}
\captura{assets/img/ethoshub-reclutador/01_dashboard/DashboardReclutador2_1.png}{Al seleccionar el botón «Vistos» nos llevara a la seccion Buscar Talento como es.}
\captura{assets/img/ethoshub-reclutador/01_dashboard/DashboardReclutador3.png}{Al observar el botón «Guardados» podemos apreciar el conteo de perfiles guardados en los últimos 7 días.}
\captura{assets/img/ethoshub-reclutador/01_dashboard/DashboardReclutador3_1.png}{Al seleccionar el botón «Guardados» podemos apreciar el conteo de perfiles guardados en los últimos 7 días como se puede ver.}
\captura{assets/img/ethoshub-reclutador/01_dashboard/DashboardReclutador4.png}{Al observar el botón «Chats» podemos apreciar el número de chats activos registrados en los últimos 7 días.}
\captura{assets/img/ethoshub-reclutador/01_dashboard/DashboardReclutador4_1.png}{Al seleccionar el botón «Chats» podemos apreciar el número de chats activos registrados en los últimos 7 días.}
\captura{assets/img/ethoshub-reclutador/01_dashboard/DashboardReclutador5.png}{Al observar el botón «Contratados» podemos apreciar el número total de candidatos contratados.}
\captura{assets/img/ethoshub-reclutador/01_dashboard/DashboardReclutador5_1.png}{Al seleccionar el botón «Contratados» podemos apreciar el número total de candidatos contratados.}
\captura{assets/img/ethoshub-reclutador/01_dashboard/DashboardReclutador6.png}{El módulo «Insight de la semana» nos muestra un resumen automático del desempeño del reclutador.}
\captura{assets/img/ethoshub-reclutador/01_dashboard/DashboardReclutador7.png}{El botón «Regenerar» nos permite actualizar el análisis de los insights de la semana.}
\captura{assets/img/ethoshub-reclutador/01_dashboard/DashboardReclutador8.png}{Después de hacer clic al botón «Regenerar» nos aparece la información actualizada en el módulo «Insight de la semana».}
\captura{assets/img/ethoshub-reclutador/01_dashboard/DashboardReclutador9.png}{El gráfico «Funnel de Conversión» nos permite visualizar el estado del proceso de los candidatos.}
\captura{assets/img/ethoshub-reclutador/01_dashboard/DashboardReclutador10.png}{El gráfico «Radar de Mercado» nos permite comparar la demanda frente a la disponibilidad de tecnologías.}
\captura{assets/img/ethoshub-reclutador/01_dashboard/DashboardReclutador11.png}{Al darle al botón de «Explorar talento» nos redirige a la sección de búsqueda de Talentos.}

\section{Motores de Búsqueda de Talento}
\label{sec:rec-busqueda}
EthosHub ofrece dos mecanismos complementarios para localizar talento: la búsqueda
inteligente en lenguaje natural (asistida por IA) y los filtros avanzados por criterios
técnicos.

\subsection{Aplicación de Filtros Avanzados y Criterios Técnicos}
\label{subsec:rec-filtros}

\captura{assets/img/ethoshub-reclutador/02_buscar-talento/BuscarTalento1.png}{Para ir al apartado de «Buscar Talento» debe hacer clic al botón «Buscar Talento».}
\captura{assets/img/ethoshub-reclutador/02_buscar-talento/BuscarTalento2.png}{Al observar el botón de flecha en la tarjeta de «Perfiles guardados» podemos apreciar la opción para ver el detalle de los candidatos guardados.}
\captura{assets/img/ethoshub-reclutador/02_buscar-talento/BuscarTalento3.png}{Al seleccionar el botón de flecha en «Perfiles guardados» nos llevara a la sección Pipelines como es.}
\captura{assets/img/ethoshub-reclutador/02_buscar-talento/BuscarTalento4.png}{Al observar el botón de flecha en la tarjeta de «Chats activos» podemos apreciar la opción para revisar las conversaciones vigentes.}
\captura{assets/img/ethoshub-reclutador/02_buscar-talento/BuscarTalento5.png}{Al seleccionar el botón de flecha en «Chats activos» nos llevara a la sección Chat como es.}
\captura{assets/img/ethoshub-reclutador/02_buscar-talento/BuscarTalento6.png}{Al observar el botón de flecha en la tarjeta de «Talentos explorados» podemos apreciar la opción para ver el historial de perfiles explorados.}
\captura{assets/img/ethoshub-reclutador/02_buscar-talento/BuscarTalento7.png}{El módulo «Búsqueda inteligente» nos permite realizar búsquedas avanzadas ingresando descripciones naturales.}
\captura{assets/img/ethoshub-reclutador/02_buscar-talento/BuscarTalento8.png}{El botón «Buscar» nos permite ejecutar la solicitud redactada en el campo de búsqueda inteligente.}
\captura{assets/img/ethoshub-reclutador/02_buscar-talento/BuscarTalento9.png}{Después de hacer clic al botón «Buscar» nos aparece el candidato compatible con los criterios ingresados.}
\captura{assets/img/ethoshub-reclutador/02_buscar-talento/BuscarTalento10.png}{Al observar el campo «TRAYECTORIA (EMPRESA)» podemos apreciar cómo la inteligencia artificial autocompletó el filtro de forma automática.}
\captura{assets/img/ethoshub-reclutador/02_buscar-talento/BuscarTalento11.png}{En caso de querer limpiar dicha busqueda seleccionamos «Limpiar todo» como podemos apreciar.}
\captura{assets/img/ethoshub-reclutador/02_buscar-talento/BuscarTalento12.png}{El botón «Desactivar match IA» nos permite limpiar los resultados actuales de la búsqueda inteligente.}
\captura{assets/img/ethoshub-reclutador/02_buscar-talento/BuscarTalento13.png}{Al observar el tablero general podemos apreciar diversos perfiles encontrados junto con el botón de «Filtros Avanzados».}
\captura{assets/img/ethoshub-reclutador/02_buscar-talento/BuscarTalento14.png}{Al observar la barra de búsqueda central podemos apreciar cómo el sistema permite realizar búsquedas manuales.}
\captura{assets/img/ethoshub-reclutador/02_buscar-talento/BuscarTalento15.png}{Después de escribir un término en la barra de búsqueda nos aparece el candidato localizado.}
\captura{assets/img/ethoshub-reclutador/02_buscar-talento/BuscarTalento16.png}{El botón «Filtros Avanzados» nos permite acceder a las opciones de segmentación.}
\captura{assets/img/ethoshub-reclutador/02_buscar-talento/BuscarTalento17.png}{Al observar el modal de filtros podemos apreciar que se ha seleccionado el filtro de nivel «Mid» y el país «Bolivia».}
\captura{assets/img/ethoshub-reclutador/02_buscar-talento/BuscarTalento18.png}{El apartado «Nivel» nos permite seleccionar entre las categorías de Junior, Mid, Senior, Lead y Architect.}
\captura{assets/img/ethoshub-reclutador/02_buscar-talento/BuscarTalento19.png}{El apartado «Experiencia» nos permite seleccionar opciones de rangos de años de trayectoria profesional.}
\captura{assets/img/ethoshub-reclutador/02_buscar-talento/BuscarTalento20.png}{El apartado «Disponibilidad» nos permite segmentar candidatos por su estado laboral.}
\captura{assets/img/ethoshub-reclutador/02_buscar-talento/BuscarTalento21.png}{La sección «TIPO DE PERFIL» nos permite elegir opciones como Frontend, Backend, DevOps, Data, Mobile, QA y Systems.}
\captura{assets/img/ethoshub-reclutador/02_buscar-talento/BuscarTalento22.png}{El campo «EDUCACIÓN (INSTITUCIÓN)» nos permite ingresar el nombre de la universidad.}
\captura{assets/img/ethoshub-reclutador/02_buscar-talento/BuscarTalento23.png}{El campo «TRAYECTORIA (EMPRESA)» nos permite especificar la empresa como criterio de búsqueda.}
\captura{assets/img/ethoshub-reclutador/02_buscar-talento/BuscarTalento24.png}{El selector de «PAÍS» nos permite elegir el país específico para la búsqueda.}
\captura{assets/img/ethoshub-reclutador/02_buscar-talento/BuscarTalento25.png}{El campo «SKILLS \& TECNOLOGÍAS» nos permite añadir nuevas habilidades mediante el botón «+».}
\captura{assets/img/ethoshub-reclutador/02_buscar-talento/BuscarTalento26.png}{El botón «Ver 1 perfiles» nos permite ejecutar la acción de filtrado.}
\captura{assets/img/ethoshub-reclutador/02_buscar-talento/BuscarTalento27.png}{Después de hacer clic al botón de filtrado nos aparece la cantidad de perfiles encontrados.}
\captura{assets/img/ethoshub-reclutador/02_buscar-talento/BuscarTalento28.png}{El botón «Limpiar todo» nos permite eliminar todos los criterios de búsqueda establecidos.}
\captura{assets/img/ethoshub-reclutador/02_buscar-talento/BuscarTalento29.png}{Después de hacer clic al botón «Limpiar todo» nos aparece una vista actualizada con los perfiles disponibles.}

\subsection{Almacenamiento de Perfiles Favoritos y Listas de Seguimiento}
\label{subsec:rec-favoritos}
Tras localizar un perfil de interés puede revisarlo en detalle, guardarlo en su
\textit{pipeline} y abrir un canal de contacto directo.

\captura{assets/img/ethoshub-reclutador/02_buscar-talento/BuscarTalento30.png}{Al seleccionar el botón «Ver Perfil» podemos apreciar el detalle completo de la trayectoria del candidato.}
\captura{assets/img/ethoshub-reclutador/02_buscar-talento/BuscarTalento31.png}{Al seleccionar el botón «Ver Perfil» podemos apreciar el detalle completo de la trayectoria del candidato.}
\captura{assets/img/ethoshub-reclutador/02_buscar-talento/BuscarTalento32.png}{Al observar el botón «Corazón» podemos apreciar la opción para guardar al candidato en la pipeline.}
\captura{assets/img/ethoshub-reclutador/02_buscar-talento/BuscarTalento33.png}{Después de hacer clic al botón «Corazón» nos aparece el mensaje «Guardado en pipeline».}
\captura{assets/img/ethoshub-reclutador/02_buscar-talento/BuscarTalento34.png}{El botón «Chat» nos permite iniciar una comunicación directa con el candidato seleccionado.}
\captura{assets/img/ethoshub-reclutador/02_buscar-talento/BuscarTalento35.png}{Al seleccionar el botón «Chat» nos redirige a la bandeja de mensajes.}

\section{Evaluación y Revisión de Portafolios}
\label{sec:rec-evaluacion}
La evaluación del talento combina la inspección de cada portafolio dentro del pipeline
de selección y la consulta de los rankings de competencia.

\subsection{Inspección de Proyectos y Evidencias Técnicas (Pipeline)}
\label{subsec:rec-pipeline}
El \textit{pipeline} organiza a los candidatos en columnas (Guardados, Revisión,
Entrevista, Oferta y Contratado) y permite compararlos técnicamente y obtener un
veredicto asistido por IA.

\captura{assets/img/ethoshub-reclutador/04_pipelines/PipelinesReclutador1.png}{Para ir al apartado de «Pipelines», debe hacer clic en el botón «Pipelines».}
\captura{assets/img/ethoshub-reclutador/04_pipelines/PipelinesReclutador2.png}{Al observar la columna de «Guardados» , podemos apreciar la lista de candidatos que han sido seleccionados previamente.}
\captura{assets/img/ethoshub-reclutador/04_pipelines/PipelinesReclutador3.png}{Al arrastrar un candidato a la columna de «Revisión», podemos avanzar su estado en el proceso de selección.}
\captura{assets/img/ethoshub-reclutador/04_pipelines/PipelinesReclutador4.png}{Al desplazar un perfil hacia la columna de «Entrevista», actualizamos su etapa en el pipeline como es.}
\captura{assets/img/ethoshub-reclutador/04_pipelines/PipelinesReclutador5.png}{Al mover un candidato a la columna de «Oferta», se registra su avance hacia la etapa final.}
\captura{assets/img/ethoshub-reclutador/04_pipelines/PipelinesReclutador6.png}{Al colocar a un perfil en la columna de «Contratado», finalizamos su flujo dentro del pipeline.}
\captura{assets/img/ethoshub-reclutador/04_pipelines/PipelinesReclutador7.png}{Al observar el botón de tres puntos en la tarjeta de un candidato , podemos apreciar las acciones adicionales disponibles para ese perfil.}
\captura{assets/img/ethoshub-reclutador/04_pipelines/PipelinesReclutador8.png}{Al seleccionar el botón de tres puntos se despliega un menú con las opciones para gestionar al candidato, como es.}
\captura{assets/img/ethoshub-reclutador/04_pipelines/PipelinesReclutador9.png}{Al hacer clic en la opción «Mover a...» del menú desplegable, podemos cambiar al candidato de etapa.}
\captura{assets/img/ethoshub-reclutador/04_pipelines/PipelinesReclutador10.png}{Al seleccionar «Mover a...» nos aparece la ventana emergente para elegir la nueva etapa del perfil.}
\captura{assets/img/ethoshub-reclutador/04_pipelines/PipelinesReclutador11.png}{Al elegir la opción «Añadir nota privada» en el menú de acciones, podemos dejar comentarios confidenciales sobre el candidato.}
\captura{assets/img/ethoshub-reclutador/04_pipelines/PipelinesReclutador12.png}{Al hacer clic en el botón «Guardar nota» después de escribir en la ventana de notas privadas, registramos la información en el perfil.}
\captura{assets/img/ethoshub-reclutador/04_pipelines/PipelinesReclutador13.png}{Al observar la sección de notas privadas dentro del perfil , podemos visualizar el historial de anotaciones guardadas.}
\captura{assets/img/ethoshub-reclutador/04_pipelines/PipelinesReclutador14.png}{Al seleccionar la opción «Abrir chat» en el menú de acciones del candidato, podemos iniciar una comunicación directa.}
\captura{assets/img/ethoshub-reclutador/04_pipelines/PipelinesReclutador15.png}{Al hacer clic a «Abrir chat» nos llevará a la sección de Chat donde aparecen las conversaciones vigentes con los candidatos, como es.}
\captura{assets/img/ethoshub-reclutador/04_pipelines/PipelinesReclutador16.png}{Al elegir la opción «Eliminar» en el menú de acciones, podemos descartar a un candidato del pipeline.}
\captura{assets/img/ethoshub-reclutador/04_pipelines/PipelinesReclutador17.png}{Después de confirmar la eliminación, aparece una notificación en la parte superior que confirma que el perfil fue removido.}
\captura{assets/img/ethoshub-reclutador/04_pipelines/PipelinesReclutador18.png}{Al marcar las casillas de selección en varios perfiles , se activa la opción para compararlos.}
\captura{assets/img/ethoshub-reclutador/04_pipelines/PipelinesReclutador19.png}{Al presionar el botón «Comparar» tras seleccionar los perfiles, se despliega la ventana de comparación como es.}
\captura{assets/img/ethoshub-reclutador/04_pipelines/PipelinesReclutador20.png}{Al observar la ventana del comparador de perfiles , podemos apreciar el desglose técnico y las habilidades de los candidatos.}
\captura{assets/img/ethoshub-reclutador/04_pipelines/PipelinesReclutador21.png}{Al visualizar la parte inferior del comparador, podemos ver el gráfico de radar técnico comparativo.}
\captura{assets/img/ethoshub-reclutador/04_pipelines/PipelinesReclutador22.png}{Al observar el área de «Veredicto Técnico IA» , podemos ingresar un rol objetivo para obtener un análisis detallado.}
\captura{assets/img/ethoshub-reclutador/04_pipelines/PipelinesReclutador23.png}{Al hacer clic en «Generar análisis» dentro del veredicto técnico, el sistema nos muestra los pros, contras y riesgos de los perfiles.}
\captura{assets/img/ethoshub-reclutador/04_pipelines/PipelinesReclutador24.png}{Al seleccionar el botón de cerrar en la ventana del comparador, volvemos a la vista principal del pipeline como es.}

\subsection{Verificación de Sellos de Validación e Historial (Rankings)}
\label{subsec:rec-rankings}
La sección \campo{Rankings} clasifica el talento por competencia y permite filtrar por
área técnica, así como conocer el porqué de cada posición.

\captura{assets/img/ethoshub-reclutador/03_rankings/RankingReclutador1.png}{Para ir al apartado de «Rankings», debe hacer clic en el botón «Rankings».}
\captura{assets/img/ethoshub-reclutador/03_rankings/RankingReclutador2.png}{Al observar el botón «General» en la parte superior , podemos apreciar la opción para visualizar el ranking de talento general.}
\captura{assets/img/ethoshub-reclutador/03_rankings/RankingReclutador3.png}{Al seleccionar el botón «General», se despliega la vista del ranking global como es.}
\captura{assets/img/ethoshub-reclutador/03_rankings/RankingReclutador4.png}{Al observar el botón «Frontend» en la parte superior , podemos apreciar la opción para filtrar el ranking por habilidades frontend.}
\captura{assets/img/ethoshub-reclutador/03_rankings/RankingReclutador5.png}{Al seleccionar el botón «Backend» en la parte superior , podemos apreciar la opción para filtrar el ranking por habilidades backend.}
\captura{assets/img/ethoshub-reclutador/03_rankings/RankingReclutador6.png}{Al observar el botón «Data» en la parte superior , podemos apreciar la opción para filtrar el ranking por habilidades en datos.}
\captura{assets/img/ethoshub-reclutador/03_rankings/RankingReclutador7.png}{Al seleccionar el botón «Infra» en la parte superior , podemos apreciar la opción para filtrar el ranking por habilidades en infraestructura.}
\captura{assets/img/ethoshub-reclutador/03_rankings/RankingReclutador8.png}{Al observar el botón «Mobile» en la parte superior , podemos apreciar la opción para filtrar el ranking por habilidades en desarrollo móvil.}
\captura{assets/img/ethoshub-reclutador/03_rankings/RankingReclutador9.png}{El botón «¿Por qué \#1?» nos permite conocer los motivos detallados por los cuales un candidato ocupa la primera posición.}
\captura{assets/img/ethoshub-reclutador/03_rankings/RankingReclutador10.png}{Al hacer clic al botón «¿Por qué \#1?» nos aparece una ventana emergente que explica la competencia del candidato seleccionado.}

\section{Gestión de Canales de Contacto Directo}
\label{sec:rec-chat}
La sección \campo{Chat} centraliza la comunicación con los candidatos, con soporte para
mensajes de texto, archivos adjuntos, notas de voz, fotografías, vídeo y emojis.

\captura{assets/img/ethoshub-reclutador/05_chat/ChatReclutador1.png}{Al hacer clic en el botón «Chat» de la barra superior, accedemos a la interfaz principal de comunicación.}
\captura{assets/img/ethoshub-reclutador/05_chat/ChatReclutador2.png}{Al seleccionar a un candidato en el panel lateral izquierdo, se despliega el historial de conversación en el área central.}
\captura{assets/img/ethoshub-reclutador/05_chat/ChatReclutador3.png}{Al redactar un mensaje en la barra inferior y presionar el icono de enviar (avión de papel), se incorpora el nuevo texto al historial.}
\captura{assets/img/ethoshub-reclutador/05_chat/ChatReclutador4.png}{Al hacer clic en el icono del clip, ubicado a la izquierda de la barra de entrada, se abre la función para adjuntar archivos.}
\captura{assets/img/ethoshub-reclutador/05_chat/ChatReclutador5.png}{Al seleccionar y cargar un archivo (como un documento .docx), este aparece con su peso en KB y permite el envío.}
\captura{assets/img/ethoshub-reclutador/05_chat/ChatReclutador6.png}{Al ver la interfaz tras adjuntar el archivo, contamos con un botón marcado con una «x» para cancelar el envío si fuera necesario.}
\captura{assets/img/ethoshub-reclutador/05_chat/ChatReclutador7.png}{Al identificar el icono del micrófono en la barra de entrada, podemos activar la herramienta para grabar notas de voz.}
\captura{assets/img/ethoshub-reclutador/05_chat/ChatReclutador8.png}{Al iniciar la grabación, se despliega una ventana emergente que muestra el tiempo transcurrido y ofrece un botón central para detener la captura, tal como aparece.}
\captura{assets/img/ethoshub-reclutador/05_chat/ChatReclutador9.png}{Durante el proceso de grabación de audio, se dispone de un icono de «x» en la esquina superior derecha de la ventana emergente para cancelar la operación.}
\captura{assets/img/ethoshub-reclutador/05_chat/ChatReclutador10.png}{Durante el proceso de grabación de audio, se dispone de un boton «reiniciar» repetir la grabacion.}
\captura{assets/img/ethoshub-reclutador/05_chat/ChatReclutador11.png}{Al observar la ventana de previsualización del audio grabado , podemos hacer clic en el botón «Enviar» para compartirlo con el candidato.}
\captura{assets/img/ethoshub-reclutador/05_chat/ChatReclutador12.png}{Al hacer clic en el icono de los tres puntos verticales sobre la barra de reproducción de audio se despliega un menú de opciones adicionales.}
\captura{assets/img/ethoshub-reclutador/05_chat/ChatReclutador13.png}{Al seleccionar el menú desplegable en la barra de audio podemos elegir entre descargar el archivo o ajustar la velocidad de reproducción.}
\captura{assets/img/ethoshub-reclutador/05_chat/ChatReclutador14.png}{Al visualizar el área de mensajes podemos confirmar que el archivo de audio ahora aparece adjunto en el chat antes de enviarlo definitivamente.}
\captura{assets/img/ethoshub-reclutador/05_chat/ChatReclutador15.png}{Al observar el historial del chat podemos ver cómo queda registrado el audio enviado junto a otros documentos previamente compartidos.}
\captura{assets/img/ethoshub-reclutador/05_chat/ChatReclutador16.png}{Al presionar el botón de disparo dentro de la herramienta de cámara capturamos la imagen deseada.}
\captura{assets/img/ethoshub-reclutador/05_chat/ChatReclutador17.png}{Al visualizar la imagen capturada tenemos la opción de hacer clic en «Repetir» si deseamos tomar una nueva fotografía.}
\captura{assets/img/ethoshub-reclutador/05_chat/ChatReclutador18.png}{Al hacer clic en el botón «Enviar» tras capturar la fotografía procedemos a compartir la imagen tomada con el candidato.}
\captura{assets/img/ethoshub-reclutador/05_chat/ChatReclutador19.png}{Al observar la esquina superior derecha del chat podemos identificar el icono de «X» destinado a eliminar la fotografia.}
\captura{assets/img/ethoshub-reclutador/05_chat/ChatReclutador20.png}{Al adjuntar un archivo como una imagen en el chat, se visualiza en la parte inferior de la interfaz con su peso en kb y la opcion de enviar la fotografia.}
\captura{assets/img/ethoshub-reclutador/05_chat/ChatReclutador21.png}{Tambien al adjuntar una fotografia en el chat, se visualiza en la parte inferior de la interfaz con una opción marcada con una «x» para cancelar el envío.}
\captura{assets/img/ethoshub-reclutador/05_chat/ChatReclutador22.png}{Al estar en el modo de grabación de video, es posible seleccionar la pestaña «Video» para configurar la herramienta antes de iniciar la captura.}
\captura{assets/img/ethoshub-reclutador/05_chat/ChatReclutador23.png}{Tras realizar una grabación de video, aparece la opción «Repetir» que permite descartar el video capturado y volver a intentarlo.}
\captura{assets/img/ethoshub-reclutador/05_chat/ChatReclutador24.png}{Al finalizar la grabación de video, se habilita el botón «Enviar» en color azul para confirmar y compartir el archivo multimedia con el candidato.}
\captura{assets/img/ethoshub-reclutador/05_chat/ChatReclutador25.png}{En la ventana de grabación de video, se puede observar un icono de «x» en la parte superior derecha que permite cerrar la herramienta sin realizar envíos.}
\captura{assets/img/ethoshub-reclutador/05_chat/ChatReclutador26.png}{Al visualizar el área de entrada de mensajes, se identifica un icono de carita feliz que permite acceder al selector de emojis.}
\captura{assets/img/ethoshub-reclutador/05_chat/ChatReclutador27.png}{Al activar el selector de emojis, se despliega un panel con categorías y una barra de búsqueda para localizar rápidamente el emoji deseado.}
\captura{assets/img/ethoshub-reclutador/05_chat/ChatReclutador28.png}{Una vez seleccionados los emojis, estos se integran en la caja de texto y se pueden enviar haciendo clic en el icono del avión de papel.}
\captura{assets/img/ethoshub-reclutador/05_chat/ChatReclutador29.png}{Al hacer clic en el icono de información «i» situado en la parte superior derecha, se abre un panel lateral con el perfil completo del candidato, incluyendo sus habilidades y datos de contacto.}
\captura{assets/img/ethoshub-reclutador/05_chat/ChatReclutador30.png}{Para eliminar un chat hacemos click en el icono de eliminacion.}
\captura{assets/img/ethoshub-reclutador/05_chat/ChatReclutador31.png}{Al confirmar la eliminación de la conversación, se muestra una ventana emergente que solicita una última validación mediante el botón «Eliminar».}
\captura{assets/img/ethoshub-reclutador/05_chat/ChatReclutador32.png}{Al completarse el proceso de eliminación, aparece una notificación en la parte superior derecha que indica que la «Conversación eliminada de tu vista», permitiendo observar la lista actualizada de otros candidatos en el panel lateral izquierdo.}

\section{Configuración de la Cuenta del Reclutador}
\label{sec:rec-config}
Desde \campo{Configuración} (icono de engranaje) se gestionan la identidad empresarial,
la personalización de la interfaz, la seguridad de la cuenta y la exportación o
eliminación de datos.

\subsection{Identidad Empresarial}
\label{subsec:rec-identidad}

\captura{assets/img/ethoshub-reclutador/06_configuracion/configuracion_reclutador_1.png}{Para ir al apartado de configuración debe hacer clic en el botón con un icono de engranaje.}
\captura{assets/img/ethoshub-reclutador/06_configuracion/configuracion_reclutador_2.png}{El botón de «Identidad» muestra la sección entera para configurar el perfil.}
\captura{assets/img/ethoshub-reclutador/06_configuracion/configuracion_reclutador_3.png}{El boton de «Seleccionar Imagen» y el campo de «pega una URL de imagen» nos permite cambiar nuestra foto de perfil.}
\captura{assets/img/ethoshub-reclutador/06_configuracion/configuracion_reclutador_4.png}{El botón de «Guardar» nos permite guardar nuestra foto de perfil.}
\captura{assets/img/ethoshub-reclutador/06_configuracion/configuracion_reclutador_5.png}{Los campos de «Nombre» y «Apellido» nos permiten colocar nuestro nombres y apellidos respectivo con un valor máximo de 80 caracteres.}
\captura{assets/img/ethoshub-reclutador/06_configuracion/configuracion_reclutador_6.png}{En el campo de «Empresa actual» podemos escribir la empresa que pertenecemos.}
\captura{assets/img/ethoshub-reclutador/06_configuracion/configuracion_reclutador_7.png}{En el campo de «Cargo en la empresa» podemos escribir nuestros cargos que tenemos en nuestra empresa.}
\captura{assets/img/ethoshub-reclutador/06_configuracion/configuracion_reclutador_8.png}{En el campo de «Ubicación» podemos escribir nuestra dirección actual con un valor máximo de 80 caracteres.}
\captura{assets/img/ethoshub-reclutador/06_configuracion/configuracion_reclutador_9.png}{El botón con icono de un marcador nos permite desplegar un mapa interactivo.}
\captura{assets/img/ethoshub-reclutador/06_configuracion/configuracion_reclutador_10.png}{Luego de hacer click en el icono de marcador vemos un modal de «Seleccionar ubicación».}
\captura{assets/img/ethoshub-reclutador/06_configuracion/configuracion_reclutador_11.png}{En el campo de «Busca una dirección o lugar...» podemos buscar un lugar mediante calles, avenida, ciudad o país.}
\captura{assets/img/ethoshub-reclutador/06_configuracion/configuracion_reclutador_12.png}{En el mapa interactivo si hacemos clic en alguna sección del mapa se muestra un icono de marcador que podemos arrastrar para marcar nuestra ubicación exacta.}
\captura{assets/img/ethoshub-reclutador/06_configuracion/configuracion_reclutador_13.png}{El botón de «Cancelar» o «X» que es nos permite al hacer clic cancelar o cerrar el modal de selección de ubicación.}
\captura{assets/img/ethoshub-reclutador/06_configuracion/configuracion_reclutador_14.png}{El botón de «Confirmar» nos permite confirmar nuestra ubicación y guardar con un autocompletado automatico.}
\captura{assets/img/ethoshub-reclutador/06_configuracion/configuracion_reclutador_15.png}{El campo de «Telefono de contacto» un telefono de contacto que tengamos con un máximo de 60 caracteres.}
\captura{assets/img/ethoshub-reclutador/06_configuracion/configuracion_reclutador_16.png}{El botón «Guardar perfil» al hacer clic guardar nuestros datos empresariales.}
\captura{assets/img/ethoshub-reclutador/06_configuracion/configuracion_reclutador_17.png}{Vemos el llenado de nuestros datos empresariales.}
\captura{assets/img/ethoshub-reclutador/06_configuracion/configuracion_reclutador_18.png}{Los datos de empresa son posibles actualizarlos una vez que modifiques los campos y le demos clic al botón «Guardar perfil».}
\captura{assets/img/ethoshub-reclutador/06_configuracion/configuracion_reclutador_19.png}{En la sección de «Sobre tu empresa» podemos ingresar una pequeña biografia empresarial con un máximo de 500 caracteres.}
\captura{assets/img/ethoshub-reclutador/06_configuracion/configuracion_reclutador_20.png}{El botón de «Guardar bio» nos permite guardar nuestra biografia empresarial.}
\captura{assets/img/ethoshub-reclutador/06_configuracion/configuracion_reclutador_21.png}{Una vez completado el campo podemos observar que se guardar correctamente.}

\subsection{Personalización de la Interfaz}
\label{subsec:rec-personalizacion}

\captura{assets/img/ethoshub-reclutador/06_configuracion/configuracion_reclutador_22.png}{Para ir a la sección de personalización debemos hacer clic en «Personalización».}
\captura{assets/img/ethoshub-reclutador/06_configuracion/configuracion_reclutador_23.png}{En la sección de «Tema de la Interfaz» podemos hacer clic en «Claro» para seleccionar un modo claro de la plataforma.}
\captura{assets/img/ethoshub-reclutador/06_configuracion/configuracion_reclutador_24.png}{El botón de «Oscuro» al hacer clic cambia a modo oscuro toda la plataforma y el botón «Sistema» se adapta al tema del Navegador.}
\captura{assets/img/ethoshub-reclutador/06_configuracion/configuracion_reclutador_25.png}{En la sección de «Idioma de la Interfaz» podemos hacer clic en «ES» «EN» «BR» para seleccionar un idioma a nuestra elección.}

\subsection{Seguridad de la Cuenta}
\label{subsec:rec-seguridad}

\captura{assets/img/ethoshub-reclutador/06_configuracion/configuracion_reclutador_26.png}{Para ir a seguridad debemos hacer clic en «Seguridad».}
\captura{assets/img/ethoshub-reclutador/06_configuracion/configuracion_reclutador_27.png}{En la sección de «Cambiar contraseña» en el botón de «Solicitar cambio de contraseña» podemos hacer clic para pedir un código en el correo asociado al momento de registrarse en la plataforma.}
\captura{assets/img/ethoshub-reclutador/06_configuracion/configuracion_reclutador_28.png}{Luego de hacer clic en la sección de «Codigo de verificación» debemos ingresar un código que nos mando la plataforma a nuestro correo asociado.}
\captura{assets/img/ethoshub-reclutador/06_configuracion/configuracion_reclutador_29.png}{Podemos ver el código que la plataforma envió a nuestro correo asociado a la cuenta.}
\captura{assets/img/ethoshub-reclutador/06_configuracion/configuracion_reclutador_30.png}{El botón de «Volver» nos permite volver y cancelar la accion de cambiar contraseña.}
\captura{assets/img/ethoshub-reclutador/06_configuracion/configuracion_reclutador_31.png}{El botón «Verificar código» nos permite verificar el código ingresado para avanzar al siguiente paso.}
\captura{assets/img/ethoshub-reclutador/06_configuracion/configuracion_reclutador_32.png}{Luego de avanzar al siguiente paso tenemos 2 campos donde debemos ingresar la nueva contraseña.}
\captura{assets/img/ethoshub-reclutador/06_configuracion/configuracion_reclutador_33.png}{Os botones con icono de ojos visibles al hacer clic nos permiten ver la contraseña que ingresamos y tambien ocultar.}
\captura{assets/img/ethoshub-reclutador/06_configuracion/configuracion_reclutador_34.png}{El botón «Actualizar contraseña» nos permite cambiar nuestra contraseña por la nueva ingresada.}
\captura{assets/img/ethoshub-reclutador/06_configuracion/configuracion_reclutador_35.png}{En la sección de «Cambiar correo electrónico» en el botón de «Solicitar cambio de correo» al hacer clic nos envia un código a nuestro correo asociado.}
\captura{assets/img/ethoshub-reclutador/06_configuracion/configuracion_reclutador_36.png}{En el campo de «Codigo de verificación» debemos ingresar el código que nos mando la plataforma al correo.}
\captura{assets/img/ethoshub-reclutador/06_configuracion/configuracion_reclutador_37.png}{Se ve el código que la plataforma nos envió para el cambio de correo electrónico.}
\captura{assets/img/ethoshub-reclutador/06_configuracion/configuracion_reclutador_38.png}{El botón «Verificar código» nos permite verificar el código ingresado para avanzar al siguiente paso.}
\captura{assets/img/ethoshub-reclutador/06_configuracion/configuracion_reclutador_39.png}{Luego de avanzar al siguiente paso tenemos un campo donde debemos ingresar un nuevo correo electrónico.}
\captura{assets/img/ethoshub-reclutador/06_configuracion/configuracion_reclutador_40.png}{El botón «Actualizar correo» al hacer clic nos permite cambiar nuestro correo por el nuevo ingresado.}
\captura{assets/img/ethoshub-reclutador/06_configuracion/configuracion_reclutador_41.png}{El botón de «Volver» nos permite volver y cancelar la accion de cambiar correo electrónico.}

\subsection{Exportación de Datos y Zona de Peligro}
\label{subsec:rec-datos}

\captura{assets/img/ethoshub-reclutador/06_configuracion/configuracion_reclutador_42.png}{En la sección de «Exportar mis datos» el botón visible «Exportar datos» al hacer clic permite que podamos descargar nuestros datos completos.}
\captura{assets/img/ethoshub-reclutador/06_configuracion/configuracion_reclutador_43.png}{Se visualiza un modal de cargar de la exportación de datos.}
\captura{assets/img/ethoshub-reclutador/06_configuracion/configuracion_reclutador_44.png}{Una vez que termine de cargar la exportación de datos nos permite elegir la ubicación donde guardaremos nuestros datos con el nombre de la plataforma y fecha.}
\captura{assets/img/ethoshub-reclutador/06_configuracion/configuracion_reclutador_45.png}{En la sección de «Zona de peligro» el botón «Eliminar cuenta» al hacer clic despliega un modal.}
\captura{assets/img/ethoshub-reclutador/06_configuracion/configuracion_reclutador_46.png}{Luego de que aparezca el modal los botones «Cancelar» y «Enviar código» visibles que al hacer clic en cancelar cancela y cierra el modal y al hacer clic en enviar código nos manda al correo asociado.}
\captura{assets/img/ethoshub-reclutador/06_configuracion/configuracion_reclutador_47.png}{Luego de pedir el codigo a nuestro correo para eliminar la cuenta vemos un modal donde nos pide el codigo de verificacion como la.}
\captura{assets/img/ethoshub-reclutador/06_configuracion/configuracion_reclutador_48.png}{Se visualiza el código que la plataforma nos mando al correo asociado.}
\captura{assets/img/ethoshub-reclutador/06_configuracion/configuracion_reclutador_49.png}{El boton de «Volver» nos permite cancelar la accion y volver hacia atras.}
\captura{assets/img/ethoshub-reclutador/06_configuracion/configuracion_reclutador_50.png}{El boton de «Confirmar eliminacion» y una vez agregado el codigo nos permite eliminar la cuenta.}

\begin{AvisoAdvertencia}
La eliminación de la cuenta es irreversible. El sistema exige un código de verificación
enviado a su correo antes de confirmar la acción.
\end{AvisoAdvertencia}

\section{Acceso desde Dispositivos Móviles (Diseño Adaptable)}
\label{sec:rec-responsive}
El panel de reclutamiento también se adapta a teléfonos y tabletas. A continuación se
muestran las vistas móviles más representativas.

\responsivePar{assets/img/ethoshub-reclutador/07_responsive/reclutador_responsive_dashboard.jpg}{Dashboard del reclutador en vista móvil.}{assets/img/ethoshub-reclutador/07_responsive/reclutador_responsive_dashboard_graficas.jpg}{Gráficas del dashboard en vista móvil.}
\responsivePar{assets/img/ethoshub-reclutador/07_responsive/reclutador_responsive_buscar_talento.jpg}{Buscar talento en vista móvil.}{assets/img/ethoshub-reclutador/07_responsive/reclutador_responsive_rankings.jpg}{Rankings en vista móvil.}
\responsivePar{assets/img/ethoshub-reclutador/07_responsive/reclutador_responsive_pipelines.jpg}{Pipelines en vista móvil.}{assets/img/ethoshub-reclutador/07_responsive/reclutador_responsive_chat_reclutador.jpg}{Bandeja de chat en vista móvil.}
\responsiveUna{assets/img/ethoshub-reclutador/07_responsive/reclutador_responsive_chat_iniciado_reclutador.jpg}{Conversación iniciada en vista móvil.}
